\documentclass[11pt]{article}
\usepackage[margin=2.3cm]{geometry}
\usepackage{amsmath,amssymb}
\usepackage{xr}
\externaldocument{main}
\usepackage[colorlinks=true,linkcolor=blue,urlcolor=blue]{hyperref}
\usepackage{enumitem}
\usepackage{microtype}
\setlength{\parskip}{0.45em}
\setlength{\parindent}{0pt}
\newcounter{rc}
\newenvironment{comment}[1]{\refstepcounter{rc}\par\medskip\noindent\textbf{Comment~#1.}\ \itshape}{\par}
\newcommand{\resp}{\par\noindent\textbf{Response.}\ }
\newcommand{\chg}{\par\noindent\textbf{Changes.}\ }

\title{Response to the reviewer\\[0.3em]\large A reliability-constrained, multi-objective framework for
minimum-mass design of steel frame structures under material and loading uncertainty}
\author{[Author name]}
\date{}

\begin{document}
\maketitle

We thank the reviewer for a careful and constructive report. We have revised the software as well as the
manuscript. Every number in the revised manuscript is regenerated from the seeded scripts in the public
repository, and the earlier results were rerun, not edited. The reviewer's comments are summarised in
italics, followed by our response and the changes made. Section, table, figure and equation numbers refer
to the revised manuscript.

\section*{Summary of the main changes}
\begin{enumerate}[leftmargin=1.5em,itemsep=0.1em]
\item A \emph{code-compliant} formulation, RBDO-C, now imposes the reliability targets \emph{and} all
deterministic AISC 360-16 strength, drift and deflection checks (Section~\ref{sec:optimisation}). It is
the main reliability-based formulation throughout the paper. The reliability-only formulation (RBDO) is
kept for comparison.
\item Every reliability estimate is reported with its failure count, an exact (Clopper--Pearson) 95\%
confidence interval for $P_f$ and $\beta$, and a three-way acceptance rule (Section~\ref{sec:mc},
Eq.~\eqref{eq:cp}).
\item A literature comparison table with eight closely related studies and a precise statement of what
is new (Section~\ref{sec:intro}, Table~\ref{tab:literature}).
\item The surrogate is validated against full second-order re-analysis for eight designs. The
validation includes 600 random samples and up to 500 targeted samples near the limit state per
design, and it gives a confidence interval for the difference in $\beta$ (Section~\ref{sec:fullmc},
Table~\ref{tab:validation}).
\item A separate treatment of model (epistemic) uncertainty with frame-wide model factors
(Section~\ref{sec:probmodel}, Table~\ref{tab:modelunc}).
\item Optimiser robustness: 2--10 seeds per formulation and frame, SAA sizes from $2\times10^4$ to
$10^5$ with three independent streams, and an exact enumeration of a reduced problem with 11\,200
designs (Section~\ref{sec:robust}).
\item Deterministic optima with drift limits $h/400$--$h/600$, section counts and an illustrative
cost index (Sections~\ref{sec:targets} and~\ref{sec:cost}).
\item A new title (``minimum-mass''), a rewritten abstract, a reproducibility statement with a
one-command reproduction script, and consistent terminology for the ULS.
\end{enumerate}

\section*{Major comments}

\begin{comment}{1 (code compliance)}
The RBDO optimum of the 6-m bay frame has an LRFD strength ratio of 1.011. State whether the
formulation includes the AISC checks, provide a code-compliant RBDO and report whether the conclusions
change.
\end{comment}
\resp The original RBDO formulation imposed only the reliability targets (ULS, drift and deflection)
and the column-continuity rule. It did not impose any deterministic check, including the LRFD strength
check. The
revised manuscript says so explicitly (Section~\ref{sec:optimisation}). We added RBDO-C, which
additionally requires every deterministic check (LRFD strength ratio $\le 1$, drift $\le h/400$,
deflection $\le L/360$). We reran it for all seven frames with at least two GA seeds each and
verified every optimum with $10^6$ independent samples (Table~\ref{tab:verification}).

The code checks are active in two of the seven frames. The 6-m bay RBDO design (LRFD ratio 1.011) is
replaced by an RBDO-C design with ratio 0.914, which is 0.32\% heavier. The low-wind RBDO design (ratio
1.002) is replaced by one that is 0.64\% heavier. In the other five frames RBDO and RBDO-C give the same
design. All RBDO-C designs satisfy every LRFD check and meet both targets under the acceptance rule. The
qualitative conclusions do not change. The mass premium of RBDO-C over DDO is $+1.13$\% to $+10.11$\%
($+2.88$\% for the base frame). It is driven by the 50-year drift reliability of the deterministic
optima ($\beta_{drift}=0.89$--$2.54$). Section~\ref{sec:discussion} explains how RBDO-C fits the code:
it is a least-mass code-compliant design that also has a stated reliability.
\chg Sections~\ref{sec:optimisation}, \ref{sec:basedesign}, \ref{sec:parametric};
Tables~\ref{tab:designs}, \ref{tab:verification}; Fig.~\ref{fig:parametric}; all figures and tables now
use RBDO-C.

\begin{comment}{2 (statistical treatment)}
Distinguish point estimates from confidence bounds, give binomial confidence intervals for $P_f$ and
$\beta$, explain acceptance of designs near the threshold, and report failure counts. The ``0.5/N lower
bound'' wording is not justified.
\end{comment}
\resp We agree. Section~\ref{sec:mc} now defines the point estimate $\hat P_f=k/N$ and the exact
Clopper--Pearson interval (Eq.~\eqref{eq:cp}), with its mapping to $\beta$. For zero failures it uses
the one-sided bound: with $N=10^6$, $P_f\le 3.0\times10^{-6}$ and $\beta\ge4.53$. The ``0.5/N''
wording has been removed.

Each verified design is reported with $N$, $k$, $\hat\beta$ and the 95\% interval
(Tables~\ref{tab:designs}, \ref{tab:targets}, \ref{tab:driftlimits}, \ref{tab:verification}). A design is
classified as \emph{above}, \emph{within} or \emph{below} the target, depending on whether the interval
lies above the target, contains it or lies below it. ``Within'' is accepted, because the target cannot
be resolved more finely at this sample size. This is stated as a convention, and its consequences are
discussed in Section~\ref{sec:discussion}. Only one final design is ``within'': the 3-storey RBDO-C
design, with $\beta_{ULS}=2.995$ [2.98, 3.01]. The others are above both targets.

We also report the stricter alternative. In the base frame, imposing the target on the one-sided lower
confidence bound inside the SAA costs 1.6\% more steel (Table~\ref{tab:saa}). Section~\ref{sec:mc}
explains the risks of common random numbers (optimisation bias and correlated estimates) and how we
control them: independent re-verification, other SAA sizes and streams, and the lower-bound variant.
\chg Section~\ref{sec:mc}; Tables~\ref{tab:designs}, \ref{tab:saa}, \ref{tab:targets},
\ref{tab:driftlimits}, \ref{tab:verification}; Figs.~\ref{fig:pareto}, \ref{fig:parametric},
\ref{fig:driftlimits}.

\begin{comment}{3 (novelty)}
Add a comparison table of 5--10 close studies and state precisely what is new.
\end{comment}
\resp Table~\ref{tab:literature} compares eight studies on reliability-based or probabilistic
optimisation of steel frames, ordered by year, with our study. Its columns are the structural model,
design variables, reliability formulation, optimiser, verification and reproducibility. The
introduction then states the contribution in four points. We do not claim novelty for any single
ingredient. The new elements are:
\begin{enumerate}[label=(\roman*),itemsep=0.1em]
\item a like-for-like comparison of DDO, DDO with tightened drift limits, RBDO and code-compliant RBDO
on the same frames, catalogue and optimiser, with interval-verified reliability;
\item serviceability reliability treated together with a system ULS;
\item a verification chain that ends in a stratified full second-order Monte Carlo validation for
several designs;
\item an open release with a one-command reproduction.
\end{enumerate}
The table entries are based on the abstracts and bibliographic records available to us; ``n.r.''
means not reported there. We have checked the bibliographic data of the six references added to the
reference list.
\chg Section~\ref{sec:intro}, Table~\ref{tab:literature}; six references added.

\begin{comment}{4 (validation breadth)}
Validate several designs (9 storeys, 10-m bays, high wind), compare failure classifications, report
the uncertainty of the difference in $P_f$ and show the accuracy near the constraints.
\end{comment}
\resp Section~\ref{sec:fullmc} now validates eight designs: the base DDO design and the RBDO-C
designs of all seven frames, including the 9-storey, 10-m bay and high-wind frames. For each design
it uses two sets of samples:
\begin{enumerate}[label=(\alph*),itemsep=0.1em]
\item 600 random samples of the wind scenario, which give the bias and scatter of the utilisation and
drift and the agreement of the drift classification;
\item a stratified tail estimator. From a pool of $2\times10^5$ samples, all samples with surrogate
utilisation $U\ge1$ (up to 250) and 250 samples with $0.85\le U<1$ are re-analysed with a full P-$\Delta$
PyNite model with member-wise random modulus and fabrication factors. This gives $P_f$ of the full
model, its ratio to the surrogate and the difference in $\beta$, with Bayesian (Jeffreys) 95\%
intervals.
\end{enumerate}
The results are in Table~\ref{tab:validation} and Fig.~\ref{fig:validation}:
\begin{itemize}[itemsep=0.1em]
\item The utilisation bias is $+0.4$ to $+1.1$\% (s.d.\ 1.0--1.6\%) and the drift bias is $-1.0$ to
$+1.4$\%.
\item 0--7 of the 600 drift classifications disagree.
\item 79--91\% of the surrogate failures also fail in the full model, and 0--4 of the 250 near misses
fail in the full model.
\item $\Delta\beta=\beta_{full}-\beta_{sur}$ lies between $-0.07$ and $+0.05$, and its interval contains
zero for seven of the eight designs.
\end{itemize}
The exception is the 3-storey RBDO-C design. There, $\Delta\beta=-0.07$ [$-0.17$, $-0.01$]; the
surrogate is slightly unconservative and the design's full-model $\beta_{ULS}$ is about 2.93. We report
this and discuss it rather than hiding it. Accuracy near the constraints is shown by the near-miss
stratum and by the largest full-to-surrogate utilisation ratio (at most 1.049).
\chg Section~\ref{sec:fullmc}; Table~\ref{tab:validation}; Figs.~\ref{fig:validation}, \ref{fig:fullmc};
new function \texttt{broad\_validation} and stage \texttt{validate}.

\begin{comment}{5 (types of uncertainty)}
Separate aleatory from epistemic and model-form uncertainty, justify the dependence assumptions, discuss
the omissions (residual stresses, imperfections, connections, inelastic redistribution) and limit the
interpretation of the system ULS.
\end{comment}
\resp
\begin{itemize}[itemsep=0.1em]
\item \emph{Uncertainty types.} Section~\ref{sec:probmodel} now distinguishes the aleatory variables
of Table~\ref{tab:rv} from model-form uncertainty. Model-form uncertainty is represented by frame-wide
lognormal factors on resistance, load effect and drift. Because the baseline indices are conditional
on a perfect model, these factors are \emph{not} part of the baseline. Their effect is quantified
separately (Table~\ref{tab:modelunc}): a c.o.v.\ of 0.05 lowers $\beta_{ULS}$ by 0.23--0.30, a c.o.v.\
of 0.10 lowers it by 0.82--0.90, and $\beta_{drift}$ changes by at most 0.05.
\item \emph{Dependence.} Section~\ref{sec:resistance} explains how the system event is built: the
union over all members and both Turkstra scenarios. Dependence enters through the shared loads,
modulus, partly correlated yield stresses ($\rho=0.5$, assumed), fabrication factor and imperfection,
and it is represented exactly by joint sampling. Section~\ref{sec:sensitivity} tests the assumed
values.
\item \emph{Omissions and interpretation.} We list the omissions (residual stresses and member
imperfections only through the AISC column curves and $\tau_b$; no inelastic redistribution; rigid,
non-failing connections). The ULS is defined consistently as ``first attainment of a member limit
state in a second-order elastic analysis, not collapse'' in the abstract, Section~\ref{sec:resistance}
and the Discussion.
\end{itemize}
\chg Sections~\ref{sec:resistance}, \ref{sec:probmodel}, \ref{sec:sensitivity},
\ref{sec:discussion}; Table~\ref{tab:rv} (new Basis/Status columns and a model-factor row);
Table~\ref{tab:modelunc}; Fig.~\ref{fig:sensitivity}.

\begin{comment}{6 (robustness of the optimiser)}
Use more seeds; report best, median and worst results; count distinct designs; use an alternative
search or reduced-space enumeration; state the stopping criterion; show the effect of the SAA size on
feasibility; do not claim global optimality.
\end{comment}
\resp
\begin{itemize}[itemsep=0.1em]
\item \emph{Seeds.} Table~\ref{tab:robust} reports best, median and worst masses, the number of
distinct final designs and the number of designs evaluated for every frame and formulation. The base
frame was run with 10 seeds (DDO) and 5 seeds (RBDO-C); the other frames with 2--3 seeds. For the base
frame, 8 of 10 DDO runs found the best design (the other two were 0.3\% heavier) and all five RBDO-C
runs found the same design. The largest spread is 1.6\%, for the high-wind RBDO-C runs, and we report
the lightest.
\item \emph{Enumeration.} A reduced problem with three variables and 11\,200 designs was solved
exactly by enumeration in order of mass (Table~\ref{tab:enum}). The GA found the exact optimum in 10 of
10 runs for each formulation, after evaluating about 700 designs.
\item \emph{SAA size.} $N=2\times10^4$, $5\times10^4$ (three independent streams) and $10^5$ all give
the same base-frame RBDO-C design (Table~\ref{tab:saa}).
\item \emph{Stopping and optimality.} The stopping criterion is stated in
Section~\ref{sec:optimisation}. We no longer claim global optimality; the reported designs are the
best found.
\end{itemize}
\chg Sections~\ref{sec:optimisation}, \ref{sec:robust}; Tables~\ref{tab:robust}, \ref{tab:enum},
\ref{tab:saa}; Fig.~\ref{fig:ga}.

\begin{comment}{7 (economics)}
Use mass as the objective and cost only as a proxy; report section counts; add a simple cost
sensitivity; examine several drift limits, not only $h/500$.
\end{comment}
\resp The objective is now called steel mass throughout, including the title. Cost is discussed only
through an illustrative cost index (Section~\ref{sec:costmodel}, Eq.~\eqref{eq:cost}). The index adds
a per-section penalty and a beam-mass term for connection work, evaluated for five parameter pairs.
Table~\ref{tab:cost} gives the number of distinct sections and the index. The ranking of DDO and
RBDO-C changes in only one case (6-m bays with a high section penalty, $-0.19$\%).

Deterministic optima were computed for drift limits $h/400$, $h/450$, $h/500$ and $h/600$ for every
frame (Table~\ref{tab:driftlimits}, Fig.~\ref{fig:driftlimits}). The lightest deterministic design
that meets both targets is $h/500$ in four frames and $h/600$ in two. In the 3-storey frame none of
these limits is enough, because that frame is strength-governed. RBDO-C is within $-5.0$\% to $+0.6$\%
of that design. It is $+0.58$\% heavier in the 6-m bay frame, because the $h/500$ design meets the
drift target only within the confidence interval and was rejected by the SAA estimate. The
manuscript reports this.
\chg Title; Sections~\ref{sec:costmodel}, \ref{sec:targets}, \ref{sec:cost};
Tables~\ref{tab:driftlimits}, \ref{tab:cost}; Fig.~\ref{fig:driftlimits}.

\section*{Additional technical issues}

\begin{comment}{T1 ($B_1$--$B_2$)}
Explain how $B_1$ and $B_2$ are recomputed and what is assumed.
\end{comment}
\resp Section~\ref{sec:superposition} now states that $B_1$ and $B_2$ are recomputed for every sample:
$P_{story}$, $P_r$ and $C_m$ follow the sampled gravity loads, and $P_{e,story}$ and $P_{e1}$ follow the
sampled modulus factor. It lists four assumptions:
\begin{enumerate}[label=(\alph*),itemsep=0.1em]
\item storey stiffness scales with the frame-wide modulus factor;
\item gravity load is distributed as in the first-order analysis;
\item there is no stiffness loss before the first member limit state;
\item the storey amplification is the same for all columns of a storey.
\end{enumerate}
Assumptions (a) and (b) are removed in the full re-analysis of Section~\ref{sec:fullmc}, and (c) is
part of the elastic ULS definition. The comparison with OpenSeesPy and PyNite P-$\Delta$
(Table~\ref{tab:opensees}, Fig.~\ref{fig:pdelta}) shows that the amplified drifts are within 1.8\% of
the second-order solution. We corrected the earlier statement that they were ``always conservative'':
the surrogate underestimates the drift by up to 1.0\%.

\begin{comment}{T2 (dependence and system event)}
Explain the dependence between scenarios and members and how the system event is built.
\end{comment}
\resp See the response to Comment~5 and Section~\ref{sec:resistance}. We also report the share of
failures by member type and scenario (Table~\ref{tab:members}).

\begin{comment}{T3 (distribution table)}
Separate sourced parameters from assumed ones.
\end{comment}
\resp Table~\ref{tab:rv} now has a ``Basis'' column with the source and a ``Status'' column (sourced /
assumed / adapted). The assumed values (bias of the yield stress 1.10, correlation 0.5,
fabrication c.o.v., model factors) are flagged and tested in Section~\ref{sec:sensitivity}.

\begin{comment}{T4 (common random numbers)}
Discuss the risks of using common random numbers.
\end{comment}
\resp Section~\ref{sec:mc}; see the response to Comment~2.

\begin{comment}{T5 (unrounded values, counts, intervals, tolerance)}
\end{comment}
\resp Percentages are now computed from unrounded masses (in kg). Indices are given to three
decimals with failure counts and intervals. The acceptance tolerance is the confidence interval
itself (Section~\ref{sec:mc}).

\begin{comment}{T6 (load paths and member-level results for more frames)}
\end{comment}
\resp Table~\ref{tab:members} reports, for all seven frames and both DDO and RBDO-C, where the failures
occur (beams, exterior or interior columns) and in which scenario. Section~\ref{sec:parametric}
explains where RBDO-C adds steel, e.g.\ $+577$\,kg in beams and $-151$\,kg in columns for 10-m bays.
Fig.~\ref{fig:memberutil} shows member utilisations of the base frame.

\begin{comment}{T7 (reproducibility)}
Give the repository URL, commit, environment and a one-command reproduction.
\end{comment}
\resp The ``Data and code availability'' statement gives the URL, folder, branch and environment
(\texttt{requirements.txt}). The script \texttt{reproduce.sh} reruns every stage and rebuilds the
manuscript. The exact commit is recorded in the repository. We will create a release tag with a
DOI on acceptance.

\section*{Terminology, figures and numerical consistency}

\begin{comment}{ULS terminology}
\end{comment}
\resp ``Component-based, second-order elastic ULS (first attainment of a member limit state, not
collapse)'' is now used consistently in the abstract, Section~\ref{sec:resistance}, the Discussion and
the Conclusions.

\begin{comment}{Fig.~2 (P-$\Delta$ comparison)}
State where the largest differences are.
\end{comment}
\resp The caption of Fig.~\ref{fig:pdelta} and the text in Section~\ref{sec:opensees} now say that
the largest differences occur in the upper storeys and that the surrogate underestimates the drift by
at most 1\%.

\begin{comment}{Fig.~3 (full Monte Carlo)}
Add classification statistics and confidence intervals.
\end{comment}
\resp These statistics are now in Table~\ref{tab:validation} and Fig.~\ref{fig:validation}. The
caption of Fig.~\ref{fig:fullmc} gives the sample sizes.

\begin{comment}{Fig.~8 (Pareto front)}
Separate verified designs from SAA estimates and show confidence intervals.
\end{comment}
\resp In Fig.~\ref{fig:pareto} the non-dominated designs are re-verified with $4\times10^5$
independent samples and shown with 95\% intervals. Designs that fail the LRFD check are drawn with
open markers. SAA-only values are drawn in grey, and the DDO and RBDO-C optima ($10^6$ samples) are
marked separately.

\begin{comment}{Numerical audit}
The 2.99 index, the 1.011 ratio, counts consistent with $\beta$, percentages from unrounded masses,
and consistent notation.
\end{comment}
\resp
\begin{itemize}[itemsep=0.1em]
\item The index 2.99 is the verified $\beta_{ULS}$ of the reliability-only (RBDO) low-wind design. It is
now reported unrounded with its count and interval: 2.986 [2.97, 3.00] ($k=1415$ in $N=10^6$). This
interval contains the target, so the design is ``within'' under the acceptance rule. The code-compliant
design for this frame has 3.019 [3.00, 3.04] (Table~\ref{tab:verification}).
\item The 1.011 ratio is explained under Comment~1.
\item The tables are generated by \texttt{make\_figures.py} directly from the result files, so counts,
indices and intervals are consistent by construction.
\item Percentages are computed from unrounded masses.
\item Notation: P-$\Delta$ denotes the storey (sway) effect and P-$\delta$ the member effect; both are
used consistently.
\end{itemize}

\section*{Minor comments}
\begin{itemize}[itemsep=0.1em]
\item \emph{Title.} Changed to ``\ldots minimum-mass design\ldots''.
\item \emph{Abstract.} Rewritten. It states the component-based second-order elastic ULS and the
model-uncertainty treatment, and it contains fewer numbers.
\item \emph{Author placeholders.} The author names and affiliations remain placeholders in this
version and will be completed before resubmission.
\item \emph{Reproducibility statement.} See T7.
\item \emph{Units and notation.} See the numerical audit above.
\item \emph{References.} Six references were added, with bibliographic data checked against publisher
or indexing records.
\item \emph{Figures and captions.} All figures were regenerated from the result files. Legends were
moved so that they do not overlap the data, and the captions state the sample sizes.
\item \emph{Repetition.} The Results and Discussion were restructured to state each finding once.
\end{itemize}

\section*{Limitations acknowledged in the revision}
\begin{itemize}[itemsep=0.1em]
\item The surrogate is slightly unconservative for the 3-storey design ($\Delta\beta=-0.07$).
\item RBDO-C is 0.58\% heavier than the $h/500$ deterministic design in the 6-m bay frame.
\item The GA spread is 1.6\% for the high-wind RBDO-C runs.
\item The literature table relies on abstracts and records.
\item The conclusions are limited to regular planar frames under gravity and wind.
\end{itemize}
These limitations are stated in the manuscript.

\end{document}
